% chapters/ch05.tex
\chapter{Modelo de diagnóstico educativo inclusivo basado en datos}

\begin{refsection}
\addbibresource{references/ch05.bib}

\section{Fundamentación del modelo}

El diagnóstico educativo constituye un componente esencial en los procesos de enseñanza y aprendizaje, particularmente en contextos educativos orientados a la inclusión. A través del diagnóstico es posible identificar características del aprendizaje, reconocer necesidades educativas diversas y orientar decisiones pedagógicas que favorezcan la equidad y la participación de todos los estudiantes.

Tradicionalmente, los diagnósticos educativos han estado basados en evaluaciones estandarizadas, pruebas psicométricas o apreciaciones docentes derivadas de la observación del desempeño académico. Si bien estos enfoques han aportado herramientas relevantes para comprender ciertos aspectos del aprendizaje, presentan limitaciones cuando se trata de analizar la complejidad de los procesos educativos en contextos caracterizados por la diversidad cultural, cognitiva y socioeducativa.

En el contexto de la sociedad del conocimiento, el creciente uso de tecnologías digitales en educación abre nuevas posibilidades para complementar los enfoques diagnósticos tradicionales mediante sistemas tecnológicos capaces de recopilar, analizar e interpretar datos educativos de manera sistemática. Estos sistemas permiten integrar múltiples fuentes de información y aplicar técnicas de analítica del aprendizaje (\textit{Learning Analytics}) orientadas a identificar patrones, generar perfiles de aprendizaje y apoyar la toma de decisiones pedagógicas basadas en evidencia.

El modelo propuesto en esta investigación se fundamenta en la articulación entre tres enfoques conceptuales desarrollados en el marco teórico (Capítulo II):

\begin{itemize}
\item la inclusión educativa como principio orientador que promueve la equidad, la participación y la eliminación de barreras para el aprendizaje;
\item el Diseño Universal para el Aprendizaje (\DUA) como marco pedagógico que propone anticipar la diversidad del alumnado mediante múltiples medios de representación, acción y compromiso;
\item la educación basada en datos y la analítica del aprendizaje como enfoques metodológicos que permiten transformar la evidencia educativa en conocimiento pedagógico útil.
\end{itemize}

A partir de esta articulación conceptual, la presente investigación propone un modelo de diagnóstico educativo inclusivo basado en datos que integra múltiples fuentes de información y utiliza herramientas de analítica del aprendizaje para apoyar la identificación de necesidades educativas y la generación de recomendaciones pedagógicas contextualizadas.

Este modelo se concibe como un sistema dinámico de producción de conocimiento pedagógico, en el cual los datos recopilados sobre los estudiantes se transforman progresivamente en diagnósticos interpretables orientados a mejorar los procesos de enseñanza y aprendizaje y a sustentar decisiones pedagógicas inclusivas.

\section{Estructura general del modelo}

El modelo de diagnóstico educativo inclusivo basado en datos se estructura como un proceso sistemático que integra la recopilación de información educativa, el análisis de datos mediante herramientas de analítica del aprendizaje y la generación de diagnósticos pedagógicos orientados a la inclusión educativa.

Desde una perspectiva conceptual, el modelo integra cuatro niveles principales:

\begin{itemize}
\item la recopilación de datos educativos provenientes de diversas fuentes;
\item el procesamiento y análisis de datos mediante técnicas de analítica del aprendizaje;
\item la generación de diagnósticos educativos inclusivos;
\item el apoyo a la toma de decisiones pedagógicas fundamentadas.
\end{itemize}

Este modelo se concibe como un proceso dinámico en el cual la información recopilada sobre los estudiantes se transforma en conocimiento pedagógico útil para los docentes y las instituciones educativas.

La Figura \ref{fig:modelo-diagnostico} presenta la estructura general del modelo propuesto.

\begin{figure}[htbp]
\centering
\resizebox{\textwidth}{!}{%
\begin{tikzpicture}[
font=\footnotesize,
node distance=14mm,
block/.style={draw, rounded corners=2mm, align=center, inner sep=8pt, text width=4cm},
center/.style={draw, rounded corners=2mm, align=center, inner sep=10pt, text width=6cm},
arrow/.style={-{Stealth[length=2.6mm]}, thick}
]

\node[block] (students)
{Datos educativos del estudiante\\
Características personales\\
Contexto educativo};

\node[block, below=of students] (instruments)
{Instrumentos diagnósticos\\
Cuestionarios educativos\\
Registros del sistema};

\node[center, below=of instruments] (analytics)
{Analítica del aprendizaje\\
Procesamiento de datos educativos};

\node[center, below=of analytics] (diagnosis)
{Motor de diagnóstico educativo inclusivo};

\node[block, below left=of diagnosis] (barriers)
{Identificación de barreras\\
para el aprendizaje};

\node[block, below right=of diagnosis] (strengths)
{Identificación de fortalezas\\
y apoyos educativos};

\node[center, below=of diagnosis] (recommend)
{Generación de recomendaciones pedagógicas\\
alineadas con \DUA};

\node[block, below=of recommend] (decisions)
{Apoyo a la toma de decisiones docentes};

\draw[arrow] (students) -- (instruments);
\draw[arrow] (instruments) -- (analytics);
\draw[arrow] (analytics) -- (diagnosis);
\draw[arrow] (diagnosis) -- (barriers);
\draw[arrow] (diagnosis) -- (strengths);
\draw[arrow] (diagnosis) -- (recommend);
\draw[arrow] (recommend) -- (decisions);

\end{tikzpicture}
}
\caption{Modelo de diagnóstico educativo inclusivo basado en datos propuesto en la investigación.}
\label{fig:modelo-diagnostico}
\end{figure}

Como se observa en la Figura \ref{fig:modelo-diagnostico}, el modelo integra diversas fuentes de información educativa y las procesa mediante herramientas de analítica del aprendizaje con el fin de generar diagnósticos pedagógicos orientados a la inclusión educativa.

Este enfoque permite superar modelos diagnósticos centrados exclusivamente en déficits, promoviendo una comprensión más amplia de las características y potencialidades del aprendizaje. Asimismo, el modelo adopta una perspectiva formativa del diagnóstico, entendiendo que el objetivo del proceso diagnóstico no es clasificar o etiquetar a los estudiantes, sino generar información útil para orientar prácticas pedagógicas inclusivas y contextualizadas.

En coherencia con el enfoque de Investigación Basada en el Diseño (\DBR) adoptado en el Capítulo III, el modelo se asume como susceptible de iteración y mejora progresiva: los resultados del diagnóstico y la retroalimentación de los actores educativos permiten refinar los instrumentos, los indicadores y las recomendaciones pedagógicas que produce el sistema.

\section{Variables del diagnóstico educativo inclusivo}

El diagnóstico educativo inclusivo requiere considerar múltiples dimensiones del aprendizaje, evitando enfoques reduccionistas que se limiten exclusivamente al rendimiento académico.

Diversas investigaciones en educación inclusiva señalan que el aprendizaje está influenciado por una combinación de factores pedagógicos, cognitivos, contextuales y socioemocionales. En consecuencia, un diagnóstico educativo inclusivo debe integrar múltiples variables que permitan comprender la complejidad del proceso educativo.

En esta investigación se propone un modelo analítico que integra cuatro grandes dimensiones de variables:

\begin{itemize}
\item variables pedagógicas;
\item variables cognitivas;
\item variables contextuales;
\item variables socioeducativas.
\end{itemize}
% --- Tabla multipágina con encabezado repetido ---
\begin{APAlongtable}
{@{}p{2.9cm}p{3.2cm}Y p{2.4cm}p{2.2cm}@{}}
{Matriz de variables del modelo de diagnóstico educativo inclusivo basado en datos.}
{tab:matriz-variables-diagnostico}

\textbf{Dimensión} &
\textbf{Variable (ejemplos)} &
\textbf{Indicadores observables (ejemplos)} &
\textbf{Instrumento / Fuente} &
\textbf{Tipo de dato} \\
\midrule
\endfirsthead

\toprule
\textbf{Dimensión} &
\textbf{Variable (ejemplos)} &
\textbf{Indicadores observables (ejemplos)} &
\textbf{Instrumento / Fuente} &
\textbf{Tipo de dato} \\
\midrule
\endhead

\midrule
\multicolumn{5}{r}{\APAtablecontinued}\\
\endfoot

\endlastfoot

\textbf{Pedagógica} &
Estrategias didácticas &
Variedad de estrategias usadas; adaptación de actividades; diversidad de recursos; nivel de diferenciación pedagógica &
Observación guiada; rúbrica docente; planificación de clase &
Cualitativo / Mixto \\

\textbf{Pedagógica} &
Participación en actividades &
Frecuencia de participación; completitud de tareas; interacción en actividades; persistencia ante tareas &
Registros del sistema; listas de cotejo; analítica de interacción &
Cuantitativo \\

\textbf{Pedagógica} &
Ajustes razonables / apoyos &
Presencia de apoyos; adecuaciones implementadas; seguimiento de apoyos; respuesta del estudiante al apoyo &
Reportes docentes; historial de apoyos; registros del sistema &
Mixto \\

\midrule

\textbf{Cognitiva} &
Ritmo de aprendizaje &
Tiempo promedio por actividad; progresión temporal; necesidad de repeticiones; tasa de finalización &
Registros del sistema; analítica del aprendizaje &
Cuantitativo \\

\textbf{Cognitiva} &
Preferencias y estilos (orientativo) &
Preferencias de representación (visual/verbal); modalidad de respuesta; patrones de interacción con recursos &
Cuestionarios; trazas digitales; autoinforme &
Mixto \\

\textbf{Cognitiva} &
Estrategias cognitivas &
Uso de planificación; autorregulación; monitoreo de comprensión; búsqueda de ayuda &
Cuestionarios; entrevistas; rúbrica de autorregulación &
Cualitativo / Mixto \\

\midrule

\textbf{Contextual} &
Recursos disponibles &
Acceso a tecnología; conectividad; disponibilidad de materiales; apoyos institucionales &
Encuesta contextual; ficha institucional; entrevista &
Mixto \\

\textbf{Contextual} &
Condiciones del entorno educativo &
Tamaño de grupo; clima de aula; barreras físicas o tecnológicas; accesibilidad de recursos &
Observación; checklist de accesibilidad; reporte institucional &
Cualitativo / Mixto \\

\textbf{Contextual} &
Variables institucionales &
Políticas inclusivas; rutas de atención; soporte profesional; prácticas de seguimiento &
Análisis documental; entrevista a directivos &
Cualitativo \\

\midrule

\textbf{Socioeducativa} &
Motivación y compromiso &
Interés sostenido; autoeficacia percibida; disposición a participar; persistencia &
Cuestionario; entrevistas; analítica de participación &
Mixto \\

\textbf{Socioeducativa} &
Factores socioemocionales &
Regulación emocional; ansiedad académica (no clínica); relación con pares; percepción de apoyo &
Cuestionarios; entrevistas; observación estructurada &
Cualitativo / Mixto \\

\textbf{Socioeducativa} &
Interacciones sociales &
Colaboración; participación en grupo; comunicación con docentes; redes de apoyo &
Observación; registros de interacción; entrevistas &
Mixto \\

\end{APAlongtable}
La Figura \ref{fig:variables-diagnostico} presenta el mapa de variables considerado en el modelo.

\begin{figure}[htbp]
\centering
\resizebox{\textwidth}{!}{%
\begin{tikzpicture}[
font=\footnotesize,
node distance=14mm,
block/.style={draw, rounded corners=2mm, align=center, inner sep=8pt, text width=4cm},
center/.style={draw, rounded corners=2mm, align=center, inner sep=10pt, text width=6cm},
arrow/.style={-{Stealth[length=2.6mm]}, thick}
]

\node[center] (diagnosis)
{\textbf{Motor de diagnóstico educativo inclusivo}\\
Integración de datos educativos};

\node[block, above left=of diagnosis] (ped)
{\textbf{Variables pedagógicas}\\
Estrategias didácticas\\
Participación en actividades};

\node[block, above right=of diagnosis] (cog)
{\textbf{Variables cognitivas}\\
Estilos de aprendizaje\\
Ritmos de aprendizaje};

\node[block, below left=of diagnosis] (context)
{\textbf{Variables contextuales}\\
Entorno educativo\\
Recursos disponibles};

\node[block, below right=of diagnosis] (social)
{\textbf{Variables socioeducativas}\\
Motivación\\
Factores socioemocionales};

\node[center, below=25mm of diagnosis] (results)
{\textbf{Resultados del diagnóstico}\\
Barreras de aprendizaje\\
Fortalezas educativas\\
Recomendaciones pedagógicas};

\draw[arrow] (ped) -- (diagnosis);
\draw[arrow] (cog) -- (diagnosis);
\draw[arrow] (context) -- (diagnosis);
\draw[arrow] (social) -- (diagnosis);
\draw[arrow] (diagnosis) -- (results);

\end{tikzpicture}
}
\caption{Mapa de variables del modelo de diagnóstico educativo inclusivo basado en datos.}
\label{fig:variables-diagnostico}
\end{figure}

Como se observa en la Figura \ref{fig:variables-diagnostico}, el modelo integra diversas dimensiones del aprendizaje.

Las variables pedagógicas permiten analizar las estrategias didácticas utilizadas en el aula, las actividades de aprendizaje propuestas y los niveles de participación del estudiante en las experiencias educativas. Las variables cognitivas aportan información relacionada con los estilos de aprendizaje, los ritmos de procesamiento de la información y las estrategias cognitivas empleadas por los estudiantes durante el proceso de aprendizaje.

Las variables contextuales consideran las condiciones institucionales, el entorno educativo y la disponibilidad de recursos pedagógicos que pueden influir en las oportunidades de aprendizaje. Por su parte, las variables socioeducativas permiten comprender factores relacionados con la motivación, los intereses personales, los aspectos socioemocionales y las interacciones sociales que forman parte del proceso educativo.

La integración de estas dimensiones permite generar diagnósticos educativos más completos y contextualizados, alineados con los principios de la educación inclusiva y con la lógica del \DUA.

\section{Funcionamiento del modelo}

El funcionamiento del modelo de diagnóstico educativo inclusivo basado en datos se desarrolla a través de varias etapas interrelacionadas que transforman los datos educativos en conocimiento pedagógico útil.

Estas etapas pueden describirse de la siguiente manera:

\begin{enumerate}

\item \textbf{Recolección de datos educativos}

En esta etapa se recopila información educativa mediante diversos instrumentos diagnósticos, cuestionarios educativos, registros de interacción con el sistema y datos contextuales relacionados con el entorno educativo.

\item \textbf{Procesamiento y preparación de datos}

Los datos recopilados son sometidos a procesos de limpieza, normalización e integración con el fin de garantizar su calidad y consistencia para el análisis posterior.

\item \textbf{Análisis mediante analítica del aprendizaje}

Se aplican técnicas de analítica del aprendizaje que permiten identificar patrones educativos relevantes, generar indicadores de aprendizaje y construir perfiles educativos de los estudiantes.

\item \textbf{Generación de diagnósticos educativos}

A partir del análisis de los datos se generan diagnósticos educativos que identifican posibles barreras para el aprendizaje, así como fortalezas y potencialidades educativas.

\item \textbf{Generación de recomendaciones pedagógicas}

Finalmente, el sistema produce recomendaciones pedagógicas alineadas con los principios del \DUA, orientadas a facilitar la personalización de las estrategias didácticas y la provisión de apoyos diferenciados.

\end{enumerate}

Este proceso permite transformar datos educativos en conocimiento pedagógico útil para apoyar la toma de decisiones docentes y promover prácticas educativas más inclusivas.

\section{Integración del modelo con el sistema tecnológico}

El modelo de diagnóstico educativo inclusivo propuesto se implementa mediante el sistema tecnológico desarrollado en el marco de la investigación. Este sistema actúa como una plataforma que permite operacionalizar el modelo conceptual a través de herramientas digitales de recopilación, análisis y visualización de datos educativos.

En particular, el sistema permite:

\begin{itemize}
\item recopilar información educativa de manera sistemática mediante instrumentos digitales;
\item almacenar y gestionar datos educativos en una base de datos estructurada;
\item analizar datos mediante herramientas de analítica del aprendizaje;
\item generar diagnósticos educativos inclusivos basados en evidencia;
\item proporcionar recomendaciones pedagógicas interpretables para los docentes.
\end{itemize}

De esta manera, el sistema tecnológico actúa como una herramienta de mediación entre los datos educativos y las decisiones pedagógicas, facilitando la aplicación práctica del modelo propuesto.

Además, la implementación del modelo mediante un sistema tecnológico permite mejorar la trazabilidad de los procesos diagnósticos, fortalecer la objetividad en la interpretación de la evidencia educativa y promover una cultura de toma de decisiones basada en datos en las instituciones educativas. Esta integración prepara el terreno para el Capítulo VI, donde se presentan los criterios de validación y la estrategia de evaluación del sistema.

\section{Aporte conceptual del modelo}

El modelo de diagnóstico educativo inclusivo basado en datos propuesto en esta investigación constituye un aporte conceptual al campo de la tecnología educativa aplicada al conocimiento.

En particular, el modelo integra de manera sistemática tres enfoques que tradicionalmente han sido abordados de forma separada en la literatura educativa:

\begin{itemize}
\item la educación inclusiva;
\item el Diseño Universal para el Aprendizaje;
\item la analítica de datos educativos.
\end{itemize}

Esta integración permite desarrollar un enfoque de diagnóstico educativo que no se limita a identificar dificultades de aprendizaje, sino que busca comprender la diversidad de los procesos educativos y generar información útil para orientar prácticas pedagógicas inclusivas.

En este sentido, el modelo constituye un marco conceptual y metodológico que puede ser adaptado a distintos contextos educativos, contribuyendo al desarrollo de sistemas tecnológicos orientados a la mejora de la equidad y la calidad educativa.

\printbibliography[heading=subbibliography,title={Referencias (Capítulo V)}]

\end{refsection}